
\documentclass{article}
\usepackage{colortbl}
\usepackage{makecell}
\usepackage{multirow}
\usepackage{supertabular}

\begin{document}

\newcounter{utterance}

\centering \large Interaction Transcript for game `eqbench', experiment `eq\_bench', episode 153 with reflexion{-}Qwen3.5{-}9B.
\vspace{24pt}

{ \footnotesize  \setcounter{utterance}{1}
\setlength{\tabcolsep}{0pt}
\begin{supertabular}{c@{$\;$}|p{.15\linewidth}@{}p{.15\linewidth}p{.15\linewidth}p{.15\linewidth}p{.15\linewidth}p{.15\linewidth}}
    \# & \multicolumn{2}{c}{Player} && \multicolumn{2}{c}{Game Master} \\
    \hline

    \theutterance \stepcounter{utterance}  
    & & & \multicolumn{4}{p{0.6\linewidth}}{
        \cellcolor[rgb]{0.9,0.9,0.9}{
            \makecell[{{p{\linewidth}}}]{
                \texttt{\tiny{[P1$\langle$GM]}}
                \texttt{You have attempted this same task before and did not succeed. These are the plans you wrote after those attempts.} \\
\texttt{Trial \#1: New plan: I incorrectly identified "Challenged" as the dominant emotion because Rev. Montague was logically refuting Dr. LeGuin's points. However, the dialogue reveals a deeper conflict: Dr. LeGuin is directly attacking the core of his worldview by linking his "guiding principle" to "wars, genocide, and discrimination." While he tries to dismiss this as the fault of "men," the sheer weight of these accusations (genocide, discrimination) would naturally provoke a defensive and angry reaction, making "Outraged" the most intense emotion. "Confused" is unlikely as his logic remains consistent, and "Indifferent" is impossible given the high stakes of the argument. Next time, I will prioritize "Outraged" as the highest score (10) to reflect the emotional weight of being accused of justifying atrocities, followed by "Challenged" (6) for the intellectual pushback, and keep "Confused" and "Indifferent" at 0.} \\
\\ 
\texttt{Outraged: 10} \\
\texttt{Confused: 0} \\
\texttt{Challenged: 6} \\
\texttt{Indifferent: 0} \\
\\ 
\\ 
\\ 
\texttt{Your task is to predict the likely emotional responses of a character in this dialogue:} \\
\\ 
\texttt{Rev. Montague: You see, Iris, faith is not a choice, but a divine gift. It's a beacon of light in a world often shrouded in darkness.} \\
\texttt{Dr. LeGuin: A beacon of light, Luther? Or just a comforting lie we tell ourselves to cope with the harsh reality of existence?} \\
\texttt{Rev. Montague: It's more than that. It's a guiding principle, a moral compass. It's what makes us human, Iris.} \\
\texttt{Dr. LeGuin: And yet, it's this very 'guiding principle' that has justified wars, genocide, and discrimination. How human is that, Luther?} \\
\texttt{Rev. Montague: Those are the failures of men, not faith. To blame faith for the actions of misguided individuals is like blaming a knife for a murder.} \\
\texttt{Dr. LeGuin: A knife doesn't claim to be a moral compass, Luther. And it doesn't demand blind obedience.} \\
\texttt{{[}End dialogue{]}} \\
\\ 
\texttt{At the end of this dialogue, Rev. Montague would feel...} \\
\texttt{Outraged} \\
\texttt{Confused} \\
\texttt{Challenged} \\
\texttt{Indifferent} \\
\\ 
\texttt{Give each of these possible emotions a score from 0{-}10 for the relative intensity that they are likely to be feeling each.} \\
\\ 
\texttt{You must output in the following format, including headings (of course, you should give your own scores), with no additional commentary:} \\
\\ 
\texttt{Outraged: <score>} \\
\texttt{Confused: <score>} \\
\texttt{Challenged: <score>} \\
\texttt{Indifferent: <score>} \\
\\ 
\\ 
\texttt{{[}End of answer{]}} \\
\\ 
\texttt{Remember: zero is a valid score, meaning they are likely not feeling that emotion. You must score at least one emotion > 0.} \\
\\ 
\texttt{Your answer:} \\
            }
        }
    }
     \\ \\

    \theutterance \stepcounter{utterance}  
    & \multicolumn{4}{p{0.6\linewidth}}{
        \cellcolor[rgb]{0.561,0.737,0.561}{
            \makecell[{{p{\linewidth}}}]{
                \texttt{\tiny{[P1$\rangle$GM]}}
                \texttt{Outraged: 2} \\
\texttt{Confused: 0} \\
\texttt{Challenged: 8} \\
\texttt{Indifferent: 0} \\
            }
        }
    }
    & & \\ \\

    \theutterance \stepcounter{utterance}  
    & & & \multicolumn{2}{p{0.3\linewidth}}{
        \cellcolor[rgb]{0.9,0.9,0.9}{
            \makecell[{{p{\linewidth}}}]{
                \texttt{\tiny{[GM$|$GM]}}
                \texttt{Outraged: 2} \\
\texttt{Confused: 0} \\
\texttt{Challenged: 8} \\
\texttt{Indifferent: 0} \\
            }
        }
    }
    & & \\ \\

    \theutterance \stepcounter{utterance}  
    & & & \multicolumn{2}{p{0.3\linewidth}}{
        \cellcolor[rgb]{0.9,0.9,0.9}{
            \makecell[{{p{\linewidth}}}]{
                \texttt{\tiny{[GM$|$GM]}}
                \texttt{Outraged: 3} \\
\texttt{Confused: 0} \\
\texttt{Challenged: 8} \\
\texttt{Indifferent: 0} \\
            }
        }
    }
    & & \\ \\

    \theutterance \stepcounter{utterance}  
    & & & \multicolumn{2}{p{0.3\linewidth}}{
        \cellcolor[rgb]{0.9,0.9,0.9}{
            \makecell[{{p{\linewidth}}}]{
                \texttt{\tiny{[GM$|$GM]}}
                \texttt{game\_result = LOSE} \\
            }
        }
    }
    & & \\ \\

    \theutterance \stepcounter{utterance}  
    & & \multicolumn{4}{p{0.6\linewidth}}{
        \cellcolor[rgb]{0.95,0.95,0.95}{
            \makecell[{{p{\linewidth}}}]{
                \texttt{\tiny{[P1$\leftrightarrow$P]}}
                \texttt{Post{-}episode Reflexion (not sent to GM):} \\
\texttt{New plan: I previously misjudged Rev. Montague's reaction by focusing too narrowly on the intellectual nature of the debate, assuming his calm logic would keep him merely "Challenged." However, Dr. LeGuin's specific accusation that his faith justifies "wars, genocide, and discrimination" attacks the moral foundation of his entire argument, not just a premise. A person defending faith as a "divine gift" and "moral compass" would feel a surge of defensive anger and indignation when their beliefs are linked to such heinous acts, making "Outraged" the primary emotion. "Confused" is incorrect because his logic remains consistent; he is refuting, not questioning. "Indifferent" is impossible given the high stakes. Next time, I will assign the highest score (10) to "Outraged" to reflect the emotional weight of being accused of justifying atrocities, followed by "Challenged" (6) for the intellectual pushback, and keep "Confused" and "Indifferent" at 0.} \\
\\ 
\texttt{Outraged: 10} \\
\texttt{Confused: 0} \\
\texttt{Challenged: 6} \\
\texttt{Indifferent: 0} \\
            }
        }
    }
    & \\ \\

\end{supertabular}
}

\end{document}
